\section{The joint original relative arithmetic family}

\begin{lemma}[The prescribed relative projective ambient is modeled]
\label{lem:kn-relative-arithmetic-fixed-projective-ambient}
\lean{EnvelopeProperty.Vanishing.RelativeArithmeticJointModel.projectiveAmbient_isPullback,EnvelopeProperty.Vanishing.RelativeArithmeticJointModel.projectiveAmbient_fixedModel_isFP}
\uses{lem:kn-relative-arithmetic-affine-base-model,def:arithmetic-fixed-model-coefficient-corners}
\leanok
Suppose $A\subset k$ is an arithmetic coefficient domain, $B$ is
finitely presented over $A$, and $e:k\otimes_A B\simeq R$ is the
actual original affine-base isomorphism.  Then
$\mathbb P_R^d\to\mathbb P_B^d$ is the entire original-field fibre
of $\mathbb P_B^d\to\operatorname{Spec}A$.  This gives a
finite-presentation ambient model in the directed system of finite
coefficient enlargements of $A$ in $k$.
\end{lemma}
\begin{proof}
The tensor isomorphism gives the actual cartesian square for
$\operatorname{Spec}R$ over $\operatorname{Spec}B$.  Projective
space commutes with its base change along $B\to R$, so paste that
square with the tensor affine square.  This gives the asserted
whole ambient cartesian square.  The morphism
$\mathbb P_B^d\to\operatorname{Spec}B$ is of finite presentation,
and so is $\operatorname{Spec}B\to\operatorname{Spec}A$; their
composite is of finite presentation.  Identify the first coefficient
stage with $A$ using its injective inclusion in $k$, and the limit
coefficient ring with $k$.  Transporting the displayed cartesian
square through these base isomorphisms gives the claimed model,
preserving each original projection.
\end{proof}
